\chapter[Band 6: Injektive Funktionen]{Band 6 auf einen Blick}
\noindent\textbf{Injektive Funktionen}\par\medskip
\noindent\emph{Lesart.} Dieser eigenständige Band fügt dem Funktionsbegriff
das Axiom der Injektivität hinzu. Für die Bild- und Urbildformeln gelten
die jeweils genannten Träger- und Teilmengenbedingungen.
\begin{BandUebersicht}
\textbf{Injektive Funktion}
\(F\colon A\inj B\) bedeutet: \(F\colon A\to B\) und
\UebFormel{x,y\in A\dsep F(x)=F(y)\vdash x=y.}
\UebQuellen{\FormulaRefAuto{Injektive Funktion}[def];
\FormulaRefAuto{Injektivität}[ax]}
&
\textit{Bildmitgliedschaft erkennt das ursprüngliche Element}
Für \(Y\subseteq A\), \(x\in A\):
\UebFormel{F(x)\in F[Y]\vdash x\in Y.}
\UebQuellen{\FormulaRefAuto{F\colon A\inj B \dsep Y\in\powerset(A)\dsep x\in A\dsep F(x)\in F[Y] \vdash x\in Y}[thm]}\\
\addlinespace[.8ex]
\textbf{Bilder unter einer Injektion}
Die Bildkonstruktion aus Band 5 wird auf eine injektive Funktion
angewandt. Sei \(F\colon M\inj N\), \(A,B\subseteq M\).
&
\textit{Schnitte werden erhalten; gleiche Bilder sind eindeutig}
\UebFormel{\begin{aligned}
&F[A\cap B]=F[A]\cap F[B],\\
&F[A]=F[B]\vdash A=B.
\end{aligned}}
\UebQuellen{\FormulaRefAuto{F\colon M\inj N \dsep A\subseteq M \dsep B\subseteq M \vdash F[A\cap B]=F[A]\cap F[B]}[thm];
\FormulaRefAuto{F\colon A\inj B \dsep X\in\powerset(A)\dsep Y\in\powerset(A)\dsep F[X]=F[Y] \vdash X=Y}[thm]}\\
\addlinespace[.8ex]
\textbf{Einschränkungen und Zielmengenwechsel}
\UebFormel{(F\restriction_C)(x)\coloneqq F(x)\quad(x\in C).}
Sei \(F\colon A\inj B\), \(C\subseteq A\), \(T\subseteq B\).
\UebQuellen{\FormulaRefAuto{RestrictionDef}[def]}
&
\textit{Injektivität bleibt erhalten}
\UebFormel{\begin{aligned}
&F\restriction_C\colon C\inj B,\\
&F\restriction_{F^{-1}[T]}\colon F^{-1}[T]\inj T,\\
&B\subseteq D\vdash F\colon A\inj D.
\end{aligned}}
\UebQuellen{\FormulaRefAuto{F\colon A\inj B\dsep C\subseteq A \vdash F\restriction_C\colon C\inj B}[thm];
\FormulaRefAuto{F\colon A\inj B\dsep T\subseteq B \vdash F\restriction_{F^{-1}[T]}\colon F^{-1}[T]\inj T}[thm];
\FormulaRefAuto{InjectiveCodomainExtension}[thm]}\\
\addlinespace[.8ex]
\textbf{Inklusionsabbildung und Komposition}
Für \(A\subseteq B\):
\UebFormel{\iota_{A,B}(x)\coloneqq x\quad(x\in A).}
Für kompatible Funktionen gilt
\UebFormel{(G\circ F)(x)\coloneqq G(F(x)).}
\UebQuellen{\FormulaRefAuto{InclusionMapDef}[def];
\FormulaRefAuto{CompositionDef}[def]}
&
\textit{Komposition von Injektionen}
\UebFormel{\begin{aligned}
&F\colon A\inj B\dsep G\colon B\inj C\\
&\hspace{12mm}\vdash G\circ F\colon A\inj C,\\
&B\subseteq C\dsep F\colon A\to B\dsep x\in A\\
&\hspace{8mm}\vdash(\iota_{B,C}\circ F)(x)=F(x).
\end{aligned}}
\UebQuellen{\FormulaRefAuto{F\colon A\inj B \dsep G\colon B\inj C\vdash G\circ F\colon A\inj C}[thm];
\FormulaRefAuto{InclusionCompositionValue}[thm]}\\
\addlinespace[.8ex]
\textbf{Fallabbildung auf disjunkten Bereichen}
Für \(A\cap B=\varnothing\), \(f\colon A\inj C\),
\(g\colon B\inj C\) verbindet \(\CaseFun{A}{B}{f}{g}\)
die beiden Zweige.
\UebQuellen{\FormulaRefAuto{CaseFunctionDef}[def]}
&
\textit{Disjunkte Bilder verhindern Kollisionen}
\UebFormel{\begin{aligned}
&f[A]\cap g[B]=\varnothing\\
&\hspace{8mm}\vdash\CaseFun{A}{B}{f}{g}\colon A\cup B\inj C.
\end{aligned}}
\UebQuellen{\FormulaRefAuto{CaseFunInjectiveDisjointImages}[thm]}\\
\addlinespace[.8ex]
\textbf{Adjunktion in einer Mengenfamilie}
\UebFormel{\begin{aligned}
\FamContaining{M}{a}&\coloneqq\{X\in M\mid a\in X\},\\
\FamAvoiding{M}{a}&\coloneqq\{X\in M\mid a\notin X\}.
\end{aligned}}
Für \(\UnionClosedFam(M)\), \(\{a\}\in M\) wird auf der
vermeidenden Teilfamilie definiert:
\UebFormel{\FranklAdj{M}{a}(X)\coloneqq X\cup\{a\}.}
\UebQuellen{\FormulaRefAuto{Enthaltende Teilfamilie}[def];
\FormulaRefAuto{Vermeidende Teilfamilie}[def];
\FormulaRefAuto{Adjunktionsabbildung einer Mengenfamilie}[def]}
&
\textit{Adjunktion ist injektiv}
Unter den links genannten Voraussetzungen:
\UebFormel{\FranklAdj{M}{a}\colon
\FamAvoiding{M}{a}\inj\FamContaining{M}{a}.}
\UebQuellen{\FormulaRefAuto{\UnionClosedFam(M)\dsep \{a\}\in M \vdash \FranklAdj{M}{a}\colon \FamAvoiding{M}{a}\inj \FamContaining{M}{a}}[thm]}\\
\addlinespace[.8ex]
\textbf{Injektiver Größenvergleich}
Für \(B\subseteq A\):
\UebFormel{\AtMostHalf{B}{A}\coloneqq
\exists F\colon B\inj A\setminus B.}
Dies definiert eine injektive Fassung von \enquote{höchstens der Hälfte}.
\UebQuellen{\FormulaRefAuto{AtMostHalf}[def]}
&
\textit{Elementarer injektiver Vergleich}
\UebFormel{b\in B\vdash\exists F\colon\{a\}\inj B.}
Eine Einermenge injiziert also in jede nichtleere Menge.
\UebQuellen{\FormulaRefAuto{SingletonInjectionIntoNonempty}[thm]}\\
\end{BandUebersicht}
